% ============================ 2. BUSINESS MODELS ==================================
\newpage
\phase{Section two}{Six ways to make money}{The same molecule supports very different businesses. They differ less in technology than in how much capital they lock up and how long they wait to be paid.}

\subsection{The six models side by side}

{\tablefont
\begin{xltabular}{\textwidth}{@{}>{\RaggedRight\arraybackslash}p{24mm} L L L L L@{}}
\caption{Business models compared. Capital figures are orders of magnitude for an entrant's first serious step.}\\
\toprule
\thead{Model} & \thead{Capital needed} & \thead{Time to profit} & \thead{Risk profile} & \thead{Exit routes} & \thead{Real examples}\\ \midrule
\endfirsthead
\toprule
\thead{Model} & \thead{Capital needed} & \thead{Time to profit} & \thead{Risk profile} & \thead{Exit routes} & \thead{Real examples}\\ \midrule
\endhead
\textbf{Developer / IPP} (each plant in its own SPV, 70–80\% debt) & ₹2,900–3,600 crore for a 100 MW green-ammonia plant with its own renewables \estd & 5–8 years: 2–3 to build, then contract cash flows & Construction, offtake and winner's-curse risk; returns fixed by contract & Sell operating SPVs to infra funds; IPO (ACME raised ₹2,900 crore in 2024 \cite{acmeipo2024} \rep) & ACME, AM Green, Avaada, NTPC REL\\
\textbf{Technology / IP licensing} & \usd{}10–100 m+ of R\&D before revenue \estd & 7–10 years & Technology risk first, then adoption risk & Licence fees; acquisition by an OEM & Versogen licensing AEM to an Indian factory \cite{enki2026}; Newtrace selling electrodes into others' stacks \cite{newtrace2026}\\
\textbf{Equipment manufacturing} & \usd{}100–400 m to reach GW scale: EH2 raised \usd{}380 m in one round \cite{eh2gv2023}; Ohmium \usd{}250 m \cite{ohmiumc2023} \rep & 5–10 years; many never get there & Overcapacity, price war with China, warranty claims & IPO or trade sale; distressed sale is common & Ohmium, Electric Hydrogen, Enapter, Plug Power\\
\textbf{Hydrogen-as-a-service} (build, own, operate on the customer's site) & ₹10–100 crore per site for 0.5–5 MW \estd & 3–5 years per site & Customer credit; utilisation; technology at small scale & Sell the portfolio to an industrial-gas major or infra fund & Hygenco for Jindal Stainless \cite{hygenco2024}; Linde and INOX Air Products in grey hydrogen\\
\textbf{Asset-light software / MRV} & ₹5–50 crore \estd & 2–4 years & Small market until rules bite; regulation can help or hurt & Acquisition by a certifier, industrial-software or carbon-market firm & Accredited verification agencies under India's 2025 scheme \cite{ghci2025}; few dedicated Indian platforms yet\\
\textbf{Aggregator / trader} & Mostly working capital and guarantees \estd & Fast if a state backs it; slow otherwise & Price and counterparty risk sits with the aggregator & Rarely sold; usually state-backed & SECI's 13 green-ammonia auctions \cite{seci2025}; Germany's H2Global, which cleared at \usd{}1,153/t in 2024 \cite{pvm2025} \rep\\
\bottomrule
\end{xltabular}}

\subsection{Scoring the models for a new Indian entrant}

\begin{table}[H]\centering\tablefont
\caption{Each model scored 1–5, where 5 is better for a capital-light newcomer.}
\begin{tabularx}{\textwidth}{@{}L c c c c c c@{}}
\toprule
\thead{Model} & \thead{Low capital} & \thead{Fast cash} & \thead{Low risk} & \thead{Margin} & \thead{Defensible} & \thead{Average}\\ \midrule
Developer / IPP & 1 & 2 & 3 & 4 & 3 & 2.6\\
Technology / IP licensing & 2 & 1 & 1 & 5 & 5 & 2.8\\
Equipment manufacturing & 1 & 1 & 1 & 2 & 2 & 1.4\\
Hydrogen-as-a-service & 3 & 3 & 3 & 4 & 3 & \textbf{3.2}\\
Asset-light software / MRV & 5 & 4 & 3 & 4 & 2 & \textbf{3.6}\\
Aggregator / trader & 3 & 2 & 2 & 2 & 2 & 2.2\\
\bottomrule
\end{tabularx}
\end{table}

\begin{finding}[Key finding \textbullet\ start light, earn the right to get heavy]
No single model is both easy to start and defensible. Software and MRV are the cheapest to start but the easiest to copy. Hydrogen-as-a-service makes real margins but needs project money per site. IP is the most defensible and the slowest. The way through is a \emph{sequence}: start where capital is lightest (data and services), use the customer relationships and operating data to win hydrogen-as-a-service contracts, and only then take on developer-scale balance-sheet risk. That sequence is also the outline of the long-term venture blueprint. \confM
\end{finding}

\begin{plainwords}[Why equipment manufacturing scores lowest]
It is the layer every government is subsidising, including ours, which is exactly why it is crowded. Global factories could make about 58 GW a year but made under 5 GW in 2025 \cite{argus2026}. When supply is twelve times demand, prices fall below cost and the weakest makers fail. A newcomer would be entering the most crowded room at the worst time.
\end{plainwords}
